% SEGMENT OLP-0022-S01
% Part: sets-functions-relations
% Chapter: functions
% Section: kinds

\documentclass[../../../include/open-logic-section]{subfiles}

\begin{document}

\olfileid{sfr}{fun}{kin}
\olsection{சார்புகளின் வகைகள்}

\begin{explain}
நாம் அடிக்கடி காணும் சில வகைச் சார்புகளுக்கு ஒரு வகைப்பாட்டை
அறிமுகப்படுத்துவது பயனுள்ளதாக இருக்கும்.

முதலில், துணைச்சார்பகத்தின் ஒவ்வோர் உறுப்பும் சார்பின் ஒரு மதிப்பாக
அமையும் பண்பைக் கொண்ட சார்புகளைக் கருதலாம். இத்தகைய சார்புகள்
!!{surjective} சார்புகள் எனப்படும். அவற்றை \olref{fig:surjective}
இல் உள்ளதுபோலப் படமாகக் காட்டலாம்.

\begin{figure}
  \olasset{assets/diagrams/surjective.tikz}
  \caption{!!^a{surjective} சார்பு, துணைச்சார்பகத்தின் ஒவ்வொரு
    !!{element}[acc]யும் ஒரு மதிப்பாகக் கொண்டிருக்கும்.}
  \ollabel{fig:surjective}
\end{figure}
\end{explain}

% SEGMENT OLP-0022-S02
\begin{defn}[!!^{surjective} சார்பு]
$f \colon A \rightarrow B$ என்ற சார்புக்கு, $f$ இன் வீச்சகமும் $B$ ஆக
இருக்கும்போது, அப்போதும் அப்போது மட்டுமே, அது
\emph{!!{surjective}} பண்புடையது எனப்படும்.
அதாவது, ஒவ்வொரு $y \in B$ க்கும் $f(x) = y$ என அமையும் குறைந்தது
ஓர் $x \in A$ இருக்க வேண்டும். குறியீடுகளில்:
\[
  (\forall y \in B)(\exists x \in A)f(x) = y.
\]
இத்தகைய சார்பை $A$ இலிருந்து $B$ க்குச் செல்லும்
!!a{surjection} என்கிறோம்.
\end{defn}

% SEGMENT OLP-0022-S03
\begin{explain}
$f$ என்பது !!a{surjection} எனக் காட்ட வேண்டுமெனில், $f$ இன்
துணைச்சார்பகத்தில் உள்ள ஒவ்வொரு பொருளும் ஏதாவது ஓர் உள்ளீடு $x$ க்கான
$f(x)$ இன் மதிப்பாக அமைகிறது எனக் காட்ட வேண்டும்.

எந்தச் சார்பும் !!a{surjection} ஐ \emph{உருவாக்குகிறது} என்பதைக்
கவனிக்கவும். ஏனெனில் $f \colon A \to B$ என்ற சார்பு தரப்பட்டால்,
$f' \colon A \to \ran{f}$ என்பதை $f'(x) = f(x)$ என வரையறுப்போம்.
$\ran{f}$ என்பது $\Setabs{f(x) \in B}{x \in A}$ எனவே
\emph{வரையறுக்கப்பட்டுள்ளதால்}, இந்தச் சார்பு $f'$ நிச்சயமாக
!!a{surjection} ஆகும்.
\end{explain}

% SEGMENT OLP-0022-S04
\begin{explain}
எந்தச் சார்பும் சாத்தியமான ஒவ்வொரு உள்ளீட்டையும் ஒரே ஒரு வெளியீட்டுடன்
இணைக்கிறது. ஆனால் வெவ்வேறு உள்ளீடுகளை ஒரே வெளியீட்டுடன் ஒருபோதும்
இணைக்காத சார்புகளும் உள்ளன. அவை !!{injective} சார்புகள் எனப்படும்.
அவற்றை \olref{fig:injective} இல் உள்ளதுபோலப் படமாகக் காட்டலாம்.
\begin{figure}
  \olasset{assets/diagrams/injective.tikz}
  \caption{!!^a{injective} சார்பு இரண்டு வெவ்வேறு உள்ளீட்டுறுப்புகளை
    ஒரே மதிப்புடன் ஒருபோதும் இணைக்காது.}
  \ollabel{fig:injective}
\end{figure}
\end{explain}

% SEGMENT OLP-0022-S05
\begin{defn}[!!^{injective} சார்பு]
ஒவ்வொரு $y \in B$ க்கும், $f(x) = y$ என அமையும் $x \in A$
அதிகபட்சம் ஒன்று மட்டுமே இருக்கும்போது, அப்போதும் அப்போது மட்டுமே,
$f \colon A \rightarrow B$ என்ற சார்பு \emph{!!{injective}}
பண்புடையது எனப்படும். இத்தகைய சார்பை $A$ இலிருந்து $B$ க்குச்
செல்லும் !!a{injection} என்கிறோம்.
\end{defn}

% SEGMENT OLP-0022-S06
\begin{explain}
$f$ என்பது !!a{injection} எனக் காட்ட வேண்டுமெனில், $f$ இன்
சார்பகத்தைச் சேர்ந்த எந்த !!{element}s ஆன $x$, $y$ க்கும்,
$f(x)=f(y)$ எனில் $x=y$ எனக் காட்ட வேண்டும்.
\end{explain}

% SEGMENT OLP-0022-S07
\begin{ex}
$f(x) = 1$ எனக் கொடுக்கப்படும் மாறிலிச் சார்பு
$f\colon \Nat \to \Nat$ என்பது !!{injective} பண்பையும்
!!{surjective} பண்பையும் கொண்டிருக்கவில்லை.

$f(x) = x$ எனக் கொடுக்கப்படும் ஒன்றாதல் சார்பு
$f\colon \Nat \to \Nat$ என்பது !!{injective} பண்பையும்
!!{surjective} பண்பையும் கொண்டுள்ளது.

$f(x) = x+1$ எனக் கொடுக்கப்படும் தொடரிச் சார்பு
$f \colon \Nat \to \Nat$ என்பது !!{injective} பண்புடையது;
ஆனால் !!{surjective} பண்புடையதல்ல.

பின்வருமாறு வரையறுக்கப்படும் $f \colon \Nat \to \Nat$ என்ற சார்பு:
\[
  f(x) =
  \begin{cases}
    \frac{x}{2} & \text{$x$ இரட்டை எனில்} \\
    \frac{x+1}{2} & \text{$x$ ஒற்றை எனில்.}
  \end{cases}
\]
!!{surjective} பண்புடையது; ஆனால் !!{injective} பண்புடையதல்ல.
\end{ex}

% SEGMENT OLP-0022-S08
\begin{explain}
பல நேரங்களில் !!{injective}, !!{surjective} ஆகிய இரண்டு
பண்புகளையும் கொண்ட சார்புகளைக் கருத விரும்புகிறோம்.
இத்தகைய சார்புகளை !!{bijective} சார்புகள் என்கிறோம்.
அவை \olref{fig:bijective} இல் காட்டப்பட்டுள்ள சார்பைப் போலிருக்கும்.
!!^{bijection}s சில சமயங்களில் \emph{ஒன்றுக்கொன்று ஒத்திசைவுகள்}
என்றும் அழைக்கப்படுகின்றன. ஏனெனில் அவை துணைச்சார்பகத்தின் உறுப்புகளைச்
சார்பகத்தின் உறுப்புகளுடன் தனித்தனியாக இணைக்கின்றன.
\begin{figure}
  \olasset{assets/diagrams/bijective.tikz}
  \caption{!!^a{bijective} சார்பு துணைச்சார்பகத்தின் உறுப்புகளைச்
    சார்பகத்தின் உறுப்புகளுடன் தனித்தனியாக இணைக்கிறது.}
  \ollabel{fig:bijective}
\end{figure}
\end{explain}

% SEGMENT OLP-0022-S09
\begin{defn}[!!^{bijection}]
$f \colon A \to B$ என்ற சார்பு !!{surjective}, !!{injective}
ஆகிய இரு பண்புகளையும் கொண்டிருக்கும்போது, அப்போதும் அப்போது மட்டுமே,
அது \emph{!!{bijective}} பண்புடையது எனப்படும்.
இத்தகைய சார்பை $A$ இலிருந்து $B$ க்குச் செல்லும் (அல்லது $A$,
$B$ ஆகியவற்றுக்கு இடையிலுள்ள) !!a{bijection} என்கிறோம்.
\end{defn}

\end{document}
